\begin{center}
    {\large \textbf{Giorno 5} \textbar\ Chill, relax first \textbar\ 13 Agosto 2024 \textbar\ \textbf{Gili Asahan}, Lombok}
\end{center}

\noindent\rule{\textwidth}{0.4pt}
\begin{figure}[h!] % Portrait images below
    \centering
    % First portrait image (on the left)
    \begin{minipage}[h]{0.48\textwidth} % Set width equal to half the page
        \centering
        \includegraphics[width=\textwidth]{images/flags/Screenshot Gili A.png}
    \end{minipage}
    \hfill
    % Text
    \begin{minipage}[h]{0.48\textwidth}
        \raggedright
        \textbf{Superficie}: 4.567 km² \\
        \vspace{0.3cm}
        \textbf{Popolazione}: 3,3 milioni \\
        \vspace{0.3cm}
        \textit{L'isola}
    \end{minipage}
\end{figure}

\landscapeLargeMargin{images/day5/PXL_20240813_023805302.MP.jpg}

\landscapeLargeMargin{images/day5/PXL_20240813_023750935.jpg}

\landscapeLargeMargin{images/day5/PXL_20240813_044505740.jpg}

\landscapeLargeMargin{images/day5/PXL_20240813_044627193.jpg}

\landscapeLargeMargin{images/day5/PXL_20240813_044723734.jpg}

\landscapeLargeMargin{images/day5/PXL_20240813_085549871.MP.jpg}

\noindent 13 Agosto 2024 \newline

La prima notte lontani dalle moschee! Peccato che mi svegli di nuovo senza motivo più volte, col risultato che quando suona la sveglia sono nuovamente in coma. Avevamo deciso di alzarci presto perché il resort offre una lezione di yoga gratuita tutte le mattine alle 7:30. Sono stanco e chiedo a Eli di rimandare, così quando poi è il momento di svegliarla davvero un'ora dopo arrivano di nuovo gli insulti. Il trucco però sta nel dirle che a una certa ora smettono di servire la colazione: qui si placa e si alza subito dal letto.

Per la colazione torniamo nella solita location dove si svolgono tutti i pasti e dal menù ordiniamo due English breakfast con uova, salsiccia, flan di riso, patate, pane tostato e fagioli. Ci servono anche un succo di ananas frullato, un piatto di frutta e ovviamente caffè in abbondanza. La frutta è buonissima, soprattutto perché la spruzziamo con il loro lime che le regala un gusto tutto nuovo.

Finita la colazione prepariamo una sacca stagna con tutto l'occorrente per trascorrere la giornata in kayak, ma non appena arriviamo al Diving Center ci viene consigliato di fare prima un po' di snorkeling nella spiaggia davanti al resort e poi di noleggiare il kayak solo per mezza giornata, così ci sarebbe costato meno. Siamo sempre più sorpresi dalla loro gentilezza, che viene messa in primo piano anche di fronte al guadagno.

Seguiamo il loro consiglio, soprattutto perché il sito davanti al resort risulta essere il migliore per lo snorkeling: lo staff è molto attento a curare e far crescere i coralli sul fondale. Indossiamo maschere e scarpette e ci immergiamo. Pesci e coralli sono meravigliosi e colorati; Elisa conosce il nome di ogni cosa che vede, ma per me si chiamano tutti pesce e sasso. L'unico che riconosco è Nemo, anche se Elisa sostiene che si chiami Pesce Pagliaccio. Sarà. Se Eli avrà voglia la lascerò descrivere dettagliatamente questa parte, altrimenti parleranno le foto.

Terminato il bagno torniamo al Diving Center e noleggiamo finalmente il kayak. Quando comunichiamo di voler fare il giro dell'isola i due ragazzi ci chiedono premurosamente se siamo sicuri di volerci avventurare così lontano, dato il vento forte. Li tranquillizziamo: abbiamo visto decisamente di peggio.

Il giro inizia ed effettivamente il vento è forte, ma il mare è tranquillo e procediamo con un buon ritmo. Il vento aiuta anche a non sentire il sole, e il risultato — dopo aver fatto snorkeling a faccia in giù con la schiena rivolta al cielo, senza aver messo prima la crema — è che me ne accorgo solo quando inizia a bruciarmi.

Arriviamo al lato nord dell'isola dove c'è una risacca e le onde formano dei cavalloni, ma ci basta pagaiare un po' al largo per evitare incidenti. Qui troviamo un altro resort che all'apparenza sembra molto bello, ma da vicino ci rendiamo conto che non si può paragonare al nostro: le casette sono tutte troppo vicine l'una all'altra, sui lettini non c'è nessuno e ricorda decisamente un villaggio vacanze, non un resort tropicale. Non ci fermiamo e proseguiamo, facendo attenzione a non rovinare i fondali: spesso raschiamo il fondo, vista l'acqua bassa anche a oltre venti metri dalla riva.

Sul lato ovest dell'isola si trova un altro centro turistico, con ristoranti e negozietti. Inizialmente pensiamo di fermarci per mangiare qualcosa, ma la calca di barche attraccate ci fa passare non solo la voglia, ma anche la fame.

Continuiamo a pagaiare verso sud costeggiando l'isola fino ad arrivare quasi al punto di partenza in meno di un'ora e mezza, e da lì attraversiamo un breve tratto di mare per raggiungere un piccolo isolotto ancora più a sud, la Gili Goleng, che misurerà sì e no duecento metri da un estremo all'altro. Qui si trova un altro ottimo spot per lo snorkeling, così attracchiamo e ci immergiamo. Questa volta i fondali sono meno affascinanti rispetto alla Coral Beach di Gili Asahan, ma nonostante ciò Elisa viene catturata dai barracuda. "Catturata" nel senso che catturano la sua attenzione, anche se un rapimento da parte di un branco di barracuda sarebbe stata una scena divertente.

La giornata sta giungendo al termine quindi riattraversiamo il lembo di mare e riportiamo il kayak al Diving Center. Qui sono nuovamente molto gentili e ci spiegano, quasi timidamente, che loro sono una struttura separata da ristorante e hotel e che quindi vanno saldati separatamente. Mi sorprende però che mi dicano di tornare più tardi o anche il giorno seguente, insistendo con "chill, relax first" e l'immancabile "take a coffee first". Sono talmente gentili che insisto per pagare subito.

A questo punto sono soddisfatto dell'uscita in mare e vorrei seguire il consiglio di chill \& relax, ma Elisa non ha visto le tartarughe. Insisto perché prima si prenda un tè con me — si è fatta l'ora della merenda gratuita al ristorante — e col cibo la convinco. Per un po'.

Prendiamo un tè fatto con i loro infusi e ci serviamo dei dolcetti cucinati da loro. Si chiamano Klepon: hanno un involucro di patata viola dolce e un interno di cioccolato fuso e burro d'arachidi. Sono una bomba e facciamo due volte il bis.

Verso le cinque del pomeriggio non riesco più a trattenere Eli, che rientra in acqua alla ricerca della tartaruga gigante che si aggira per quelle zone. Ormai il sole è basso e il fondale è buio, quindi non osa allontanarsi: la vedo in piedi, dove si tocca, con solo la testa immersa con la maschera in un fondale alto sì e no un metro. Vista da lontano sembra uno struzzo, ma perdo l'attimo e non riesco a fotografarla.

Improvvisamente mi sento chiamare a gran voce. Penso che abbia trovato una tartaruga che le gira tra le dita dei piedi, ma scopro che in realtà vuole che mi butti in acqua con lei per accompagnarla nei fondali bui. Se fossi stato stanco come i giorni precedenti probabilmente mi avrebbe preso il solito attacco di ridarola, ma stavolta esprimo solo tutta la mia disapprovazione: un modo carino per dire che ho risposto con uno "stocazzo" e non sono andato.

Eli esce sconfitta e infreddolita e sono pronto ad accoglierla con l'asciugamano, e così la riconquisto.

Torniamo in camera a prepararci per la nostra ultima cena in quel paradiso tropicale. Non mi sento mai del tutto pulito perché l'acqua della doccia e del lavandino è leggermente salata — pensiamo si tratti di acqua di mare depurata. Se ci aggiungo che mi piace come mi stanno i capelli dopo che si sono asciugati dall'acqua salata e che quindi decido di non lavarli, mi sembra di rivivere i giorni delle vacanze in kayak. Al ristorante per fortuna nessuno lo nota, come nessuno nota la camicia di lino bianco che indosso da tre giorni, a completare lo stile les vagabonds.

Abbiamo saltato il pranzo e questa sarà l'ultima cena, quindi decidiamo di strafare. Io ordino per antipasto gyoza di carne fatti a mano, per primo un Beef Rendang — carne stile goulash con riso e verdure, piatto tipico indonesiano — e per dolce lo strudel della casa. Elisa invece ordina spring rolls di verdure, Tomato Sambal Tiger Prawn — nome molto chic per gamberi flambé con pomodorini abbrustoliti — ed edamame indonesiani per dolce. Da bere acqua per me e per Eli un tè fresco a base di zenzero a pezzi e miele indonesiano, dissetante e antiossidante, a detta sua.

Stavolta mi va un po' di sfiga col dolce: dapprima non hanno lo strudel, che sostituisco con la cheesecake, poi si dimenticano di portarmela e devo andare a chiederla al bancone. Errore da poco: si erano confusi col ragazzo del tavolo a fianco che aveva ordinato lo stesso dolce. Anche a Eli non va benissimo perché hanno finito gli edamame, così ordina la zuppa di mais per dessert, dal momento che aveva passato tutta la cena a dire di volerla provare. L'ha divorata, ma a me sembrava mangime per galline caldo e tritato.

I camerieri sono gentilissimi: si scusano ogni volta che arrivano al tavolo, ascoltano diligentemente l'ordine e infine chiedono se possono ripeterlo per verificare di aver capito tutto. Anche per la dimenticanza del dolce non perdono punti.

Tornati in camera finiamo di preparare le valigie, io scrivo un po' di diario mentre Eli legge i precedenti tra le lacrime dal ridere, e impostiamo la sveglia alle sette, decisi questa volta ad andare a fare yoga.

Quando è ora di andare a letto Eli mi chiede di cambiare lato perché solo dal mio c'è la lampada e la presa per il caricabatterie. Acconsento mostrando grande magnanimità, nascondendole però che dal mio lato ci sono le molle sfondate che si piantano nella schiena, e che nel pomeriggio mi ci ero coricato prima di docciarmi riempiendolo di sabbia. Quando realizza che brutto affare ha fatto è troppo tardi: mi metto la mascherina sugli occhi e mi addormento ridendo.


\pagestyle{empty} % disable numbers

%coppia
\twoImagesLargeMargins{images/day5/PXL_20240813_011113398.jpg}
                      {images/day5/PXL_20240813_011755110.jpg}
%

%coppia
\splitImageLeft{images/day5/PXL_20240813_023552231.MP.jpg}
\splitImageRight{images/day5/PXL_20240813_023552231.MP.jpg}
%

%coppia
\landscapeThinMargin{images/day5/PXL_20240813_023625263.jpg}
\twoImagesLargeMargins{images/day5/PXL_20240813_024001308.MP.jpg}
                      {images/day5/PXL_20240813_024158155.PORTRAIT.jpg}
%

%coppia
\twoImagesLargeMargins{images/day5/PXL_20240813_024201606.PORTRAIT.jpg}
                      {images/day5/PXL_20240813_024244771.PORTRAIT.jpg}
\landscapeThinMargin{images/day5/PXL_20240813_024356005.PORTRAIT.jpg}

%coppia 
\landscapeThinMargin{images/day5/PXL_20240813_024410490.PORTRAIT.jpg}
\landscapeThinMarginCentered{images/day5/PXL_20240813_023606568.MP.jpg}
%

%coppia
\splitImageLeft{images/day5/PXL_20240813_062959408.MP.jpg}
\splitImageRight{images/day5/PXL_20240813_062959408.MP.jpg}
%

%
\eightLandscapeImages{images/Kodak/100_0721.JPG}
                     {images/Kodak/100_0731.JPG}
                     {images/Kodak/100_0736.JPG}
                     {images/Kodak/100_0751.JPG}
                     {images/Kodak/100_0755.JPG}
                     {images/Kodak/100_0766.JPG}
                     {images/Kodak/100_0918.JPG}
                     {images/Kodak/100_0932.JPG}
\eightLandscapeImages{images/Kodak/100_0916.JPG}
                     {images/Kodak/100_0917.JPG}
                     {images/Kodak/100_0918.JPG}
                     {images/Kodak/100_0919.JPG}
                     {images/Kodak/100_0922.JPG}
                     {images/Kodak/100_0923.JPG}
                     {images/Kodak/100_0926.JPG}
                     {images/Kodak/100_0950.JPG}
%

%coppia
\splitImageLeft{images/Kodak/100_1001.JPG}
\splitImageRight{images/Kodak/100_1001.JPG}
%

%coppia
\twoImagesLargeMargins{images/day5/PXL_20240813_082812202.MP.jpg}
                      {images/day5/PXL_20240813_082850088.MP.jpg}
\landscapeThinMargin{images/day5/PXL_20240813_085410252.MP.jpg}
%

%coppia
\landscapeThinMargin{images/day5/PXL_20240813_085524636.jpg}
\landscapeThinMargin{images/day5/PXL_20240813_085441683.MP.jpg}

%

%coppia
\landscapeThinMargin{images/day5/PXL_20240813_085818819.jpg}
\landscapeThinMargin{images/day5/PXL_20240813_090155642.MP.jpg}
%

%coppia
\splitImageLeft{images/day5/PXL_20240813_090125354.MP.jpg}
\splitImageRight{images/day5/PXL_20240813_090125354.MP.jpg}
%

%coppia
\landscapeThinMargin{images/day5/PXL_20240813_085607769.MP.jpg}
\sixLandscapeImages{images/day5/PXL_20240813_122006791.jpg}
                   {images/day5/PXL_20240813_122010285.MP.jpg}
                   {images/day5/PXL_20240813_123335393.jpg}
                   {images/day5/PXL_20240813_123339071.MP.jpg}
                   {images/day5/PXL_20240813_131948249.MP.jpg}
                   {images/day5/PXL_20240813_133711401.jpg}
%

%coppia
\splitImageLeft{images/day5/PXL_20240813_135055536.NIGHT.jpg}
\splitImageRight{images/day5/PXL_20240813_135055536.NIGHT.jpg}
%

%%% TODO: capire se serve questa immagine
\landscapeLargeMargin{images/day5/PXL_20240812_045941652.jpg}

\pagestyle{fancy} % enable numbers